\chapter{Discusión}
\label{cap:discusion}

Este capítulo responde las preguntas de investigación, compara los resultados con
trabajos relacionados y discute los hallazgos inesperados y sus implicaciones
prácticas.

\section{Respuesta a las preguntas de investigación}

\textbf{RQ1 (rendimiento).} En el entorno local sin throttling, el acceso
híbrido CRUD/ORM + stored procedures presenta una latencia media de
23.01\,ms (IC\,95\,\% [$-7.88$; $53.91$]\,ms, n = 3) con 0\,\% de errores y
p95 máximo de 173.13\,ms, cumpliendo el umbral de 200\,ms (serie K1--K3,
2026-07-29, anterior a la implementación real de \texttt{@Cacheable} en
este endpoint: no mide un efecto de caché, solo la estabilización tras el
arranque en frío del proceso). La media observada es consistente con los
resultados de pruebas de carga de sistemas web reportados en la
literatura \cite{jiang2015,jiang2009}. El efecto de la caché de
aplicación, medido después con metodología válida (K9, K10--K14, una vez
implementado \texttt{@Cacheable}), es real en la corrida piloto (mediana
514.75\,ms fría vs.\ 191.59\,ms caliente, 63\,\% menos) pero no alcanza
significancia estadística en las cinco corridas de contraste
($p=0.6015$, $d$ de Cliff $=-0.20$): la condición caliente resultó igual
o más lenta que la fría en tres de cinco corridas, así que con $n=5$ no
se puede distinguir el efecto de caché del ruido de red del plan
gratuito de Render.

\textbf{RQ2 (seguridad).} El sistema mitiga A01, A02, A03, A05, A07 y A09 del
OWASP Top 10 \cite{owasp2021} mediante controles verificables: autenticación JWT
con cookie HttpOnly, BCrypt, repositorios Spring Data (sin SQL concatenado) y
limitación de tasa. Los riesgos identificados en la literatura de microservicios
\cite{soldani2018,waseem2021,taibi2018} se abordan con \texttt{@Procedure} y
\texttt{@Query} constantes.

\textbf{RQ3 (usabilidad).} La media SUS de 63,0 (IC\,95\,\% [53,07; 72,93], n=10)
clasifica al sistema como \emph{OK} en la escala adjetiva, dentro de la zona
marginal \cite{bangor2009,bangor2008}. Comparado con productos de software
generalistas medidos en \cite{kortum2013}, el valor queda por debajo del rango
típico de aplicaciones web ($\approx$67--74); el IC 95\% incluye valores por debajo de 70, por
lo que se recomienda iterar sobre la complejidad percibida de los ítems pares.
Se recolectaron 5 respuestas adicionales (P11--P15) que se excluyeron del
análisis oficial por falta de respaldo documental verificable de la fecha de
recolección (ver \texttt{dataset/sus/PARTICIPANTES-NO-INCLUIDOS.md}).

\textbf{RQ4 (reproducibilidad y calidad web).} El proyecto es reproducible: todos
los artefactos (datos crudos, scripts, figuras) están versionados y el protocolo
FAIR se verifica en \texttt{docs/checklists/fair.md}. Lighthouse reporta escritorio
95/91/92/90, alcanzando los cuatro umbrales; en móvil 77/91/92/90, con Performance
por debajo del umbral de 80, en línea con SPA modernos bien configurados.

\section{Comparación con trabajos relacionados}

- \emph{Rendimiento:} a diferencia de estudios que reportan percentiles sin
  intervalo de confianza \cite{jiang2009}, aquí se incluye IC\,95\,\% con t de
  Student sobre corridas independientes, siguiendo la guía de experimentación en
  ingeniería de software \cite{wohlin2012}.
- \emph{Seguridad:} el enfoque de auditoría estática (\texttt{audit-sql-dynamic.sh}
  y SpotBugs) complementa la evidencia dinámica OWASP de \cite{owasp2021},
  emulando prácticas de revisión automatizada propuestas en la literatura de
  calidad de microservicios \cite{taibi2018}.
- \emph{Usabilidad:} el uso de SUS con 10 participantes y la escala adjetiva sigue
  el protocolo estándar \cite{brooke1996,bangor2009}; la discusión sobre tamaño
  muestral mínimo y factor estructural de la escala se respalda en \cite{lewis2009}.
- \emph{Cobertura:} se reconocen las limitaciones de la cobertura como predictor
  de efectividad \cite{inozemtseva2014}, por lo que se combina con pruebas de
  integración y análisis de ramas \cite{gopinath2014,gligoric2013,ivankovic2019}.

\section{Hallazgos inesperados}

1. La primera corrida k6 concentra el p95 más alto (173.13\,ms) mientras las dos
   posteriores estabilizan p95 por debajo de 40\,ms; interpretado como arranque del
   proceso y del pool de conexiones HikariCP, no como regresión. Esta serie
   (K1--K3, 2026-07-29) es anterior a la implementación real de
   \texttt{@Cacheable} en este endpoint (commit \texttt{3f78036},
   2026-09-17), así que la bajada no puede atribuirse al caché de Redis
   -- ver \S\ref{sec:perf}.
2. El SEO de Lighthouse (90) quedó al límite pese a un frontend SPA sencillo; la
   causa es la ausencia de metadatos \emph{meta tags} completos en el \texttt{index.html}.
3. La cobertura final (99.7\,\% líneas) resultó muy superior al umbral inicial
   (14.69\,\%), gracias a la migración a JaCoCo 0.8.14 compatible con JDK 25.

\section{Implicaciones prácticas}

- Para la operación real en cooperativas, el punto de mejora inmediato es la
  usabilidad (SUS 63,0): simplificar los formularios de alta/edición.
- Para el operador de plataforma, la bajada de p95 de 173 a $<$40\,ms entre la
  primera corrida local y las dos siguientes es arranque de proceso/conexión,
  no caché (ver el hallazgo anterior); el efecto real del caché de Redis
  (K9, K10--K14) es de apenas 15--35\,ms por corrida y no alcanza
  significancia estadística con $n=5$ (\S\ref{sec:perf}), así que
  dimensionar la política TTL a partir de esta evidencia sería prematuro --
  hace falta ampliar el número de corridas antes de recomendar un valor.
- La reproducibilidad del artefacto (FAIR) simplifica la lectura y replicación del
  estudio, según exigen los estándares empíricos \cite{ralph2021,prisma2021}.